\chapter{Distributed Compute and Schedulers}\label{ch:pl:distributed-compute-and-schedulers}

A \omterm{def:pl:distributed-compute-and-schedulers:sweep}{parameter sweep} of ten thousand backtests was submitted at 18:00 on a Friday and finished at 13:47 on the Saturday. From 10:20 that morning,
fewer than a tenth of the cluster's cores had been busy. Forty of the tasks were thirty times longer than the median, and they had
been submitted last: the cluster reached them after twelve hours and then waited seven and a half for them. Nothing had failed; the scheduler had done
exactly what it was told. This chapter is about what to tell it: how a sweep becomes a \omterm{def:pl:distributed-compute-and-schedulers:graph}{task graph}, how schedulers order the graph on a
cluster, why the last few tasks decide when a sweep ends, and what caching, retries and backups cost and save.

\section{Clusters and what they are for}

\begin{definition}[Compute cluster, job scheduler]\label{def:pl:distributed-compute-and-schedulers:cluster}
A \emph{compute cluster}\index{compute cluster} is a set of machines (nodes), each with cores, memory and sometimes accelerators, managed
as one pool of capacity. A \emph{job scheduler}\index{job scheduler} decides which submitted work runs on which part of the pool and when,
according to a policy, the resources each piece of work requests and the priorities of the users who submitted it.
\end{definition}

A research platform uses its cluster for work that is large, repetitive and mostly independent: \omterm{def:pl:distributed-compute-and-schedulers:sweep}{parameter sweeps}, walk-forward refits
(Book~7, chapter~20), simulations over seeds (chapters~11 and~12), risk scenarios (chapters~18 and~20), model training (Book~12). Most of it
is the easiest kind of parallel work there is.

\begin{definition}[Embarrassingly parallel workload, parameter sweep]\label{def:pl:distributed-compute-and-schedulers:sweep}
A workload is \emph{embarrassingly parallel}\index{embarrassingly parallel workload} when it splits into tasks that do not communicate
while they run. A \emph{parameter sweep}\index{parameter sweep} is such a workload: the same computation, typically a backtest, run once
for each point of a grid of parameters, data slices or seeds.
\end{definition}

\omterm{def:pl:distributed-compute-and-schedulers:sweep}{Embarrassingly parallel} does not mean embarrassingly easy to finish on time. The tasks of a sweep have unequal durations --- a longer
lookback, a busier market, a slower data slice --- and the sweep ends only when its last task does.

\begin{definition}[Makespan, straggler]\label{def:pl:distributed-compute-and-schedulers:makespan}
The \emph{makespan}\index{makespan} of a set of tasks is the time from the start of the first to the end of the last. A
\emph{straggler}\index{straggler} is a task that finishes long after most of its peers, because it is intrinsically long, because it runs
on a slow or overloaded machine, or because it failed and was run again.
\end{definition}

\section{Schedulers and queues}

The simplest scheduler keeps one queue and gives the next free core the task at its head: first in, first out. Graham (1969) analysed this
\emph{list scheduling} and bounded it: whatever the order, the \omterm{def:pl:distributed-compute-and-schedulers:makespan}{makespan} on $m$ identical cores is at most $2 - 1/m$ times the best possible,
and taking the tasks longest first --- \emph{longest processing time}, LPT --- tightens the bound to $4/3 - 1/(3m)$. The bounds rest on a
simple observation, which is also the chapter's measuring stick.

\begin{proposition}[Lower bound on the makespan]\label{prop:pl:distributed-compute-and-schedulers:bound}
No schedule of tasks with durations $d_i$ on $m$ cores finishes before $\max\bigl(\sum_i d_i / m,\ \max_i d_i,\ L\bigr)$, where $L$ is the
longest chain of dependent tasks: the work must be done by $m$ cores, the longest task runs on one core, and a chain runs in sequence.
\end{proposition}

\omcode{code/firm/jobgraph/firm_jobgraph.py}{279}{292}{The lower bound of \cref{prop:pl:distributed-compute-and-schedulers:bound}, with the
longest chain computed over the task graph.}\label{lst:pl:distributed-compute-and-schedulers:bound}

LPT needs to know which tasks are long. A scheduler does not know durations, only estimates --- from the parameters (a longer window, a
longer history), from earlier runs of the same code, or from the user. The chapter's estimates are right to within a factor $e^{0.3z}$ with
$z$ standard normal, a typical spread for a model of run time.

\begin{definition}[Fair-share scheduling]\label{def:pl:distributed-compute-and-schedulers:fairshare}
\emph{Fair-share scheduling}\index{fair-share scheduling} gives each user or team a target share of the cluster and orders the queue so that
those using less than their share are served first; a team that submits ten thousand tasks does not lock out a team that submits five hundred
an hour later.
\end{definition}

\omcode{code/firm/jobgraph/firm_jobgraph.py}{109}{129}{Fair share in its simplest form: the next task comes from the team using the fewest
cores per unit of its target.}\label{lst:pl:distributed-compute-and-schedulers:fairshare}

\begin{dated}{2026-09}{Fair share in a production scheduler}\label{dat:pl:distributed-compute-and-schedulers:slurm}
Slurm, a widely used open-source workload manager, computes a fair-share factor from normalised shares $S$ and normalised usage $U$
(its ``Fair Tree'' algorithm, the default since Slurm 19.05): an association's level fair-share is $S/U$, above 1 when it is under-served,
and users are ranked so that the children of a better-served account always rank below those of a worse-served one.
\end{dated}

\begin{definition}[Work stealing]\label{def:pl:distributed-compute-and-schedulers:stealing}
\emph{Work stealing}\index{work stealing} is a decentralised scheduling policy: each worker keeps its own queue of tasks and takes from its
head; a worker whose queue is empty takes a task from the far end of another worker's queue. Blumofe and Leiserson (1999) analysed it
for multithreaded computations; it balances load without a central queue.
\end{definition}

\omcode{code/firm/jobgraph/firm_jobgraph.py}{132}{154}{Work stealing: tasks dealt round robin to per-core deques, an idle core stealing
from the fullest.}\label{lst:pl:distributed-compute-and-schedulers:stealing}

\section{Task graphs}

\begin{definition}[Task graph]\label{def:pl:distributed-compute-and-schedulers:graph}
A \emph{task graph}\index{task graph} is a directed acyclic graph whose nodes are tasks, each with a resource request and an estimated
duration, and whose edges say which task's output another task reads; a task is ready when all its predecessors have finished.
\end{definition}

A sweep is rarely a flat list. Its backtests read features; the features read cleaned data; the data were built once for many sweeps.
Written as a graph, each shared stage is one task that many backtests depend on, and its output can be cached under the content key of
Book~7's research workflow (chapter~29): the hash of its code, parameters and inputs. The next sweep over the same data finds the features
in the cache and runs only the backtests.

\section{Parameter sweeps and stragglers}

\begin{omfigure}
\begin{tikzpicture}[font=\scriptsize,>=Stealth]
\foreach \i/\l in {0/{feature stage 1},1/{feature stage 2},2/{\dots},3/{feature stage 100}} {
  \node[draw,rounded corners=2pt,fill=omMeth!15,minimum width=2.2cm] (f\i) at (\i*3,1.4) {\l};}
\foreach \i in {0,1,3} {\foreach \j in {-1,0,1} {\draw[->] (f\i.south) -- ++(\j*0.8,-0.8);}}
\foreach \i in {0,1,3} {\node[font=\tiny] at (\i*3,0.1) {100 backtests};}
\node[draw,rounded corners=2pt,fill=gray!10,align=center,minimum width=10cm] (c) at (4.5,-0.6) {cluster: 24 nodes $\times$ 16 cores;
one node four times slower; running work fails at 0.002 per core-hour};
\end{tikzpicture}

\omcaption{The chapter's \omterm{def:pl:distributed-compute-and-schedulers:graph}{task graph} and cluster. A hundred feature stages of twenty minutes each, each read by a hundred backtests; every
backtest's duration drawn from the measured spread, forty of them thirty times the median in the flat sweep. The feature stages are cached
by content key.}\label{fig:pl:distributed-compute-and-schedulers:graph}
\end{omfigure}

The chapter's cluster is simulated (\texttt{firm.jobgraph}); its tasks are not invented. Twenty-four small backtests of chapter~11 ---
level-3 replays of five-minute \texttt{firm.tape} sessions --- were run on this laptop through the same scheduler interface, with one
worker: their median was 0.089~seconds, the spread of their log durations 0.76, and the longest took 8.0 times the median. The simulated
sweep keeps that spread and scales the median to fifteen minutes: 10\,000 backtests, forty of them thirty times the median (seven and a
half hours, a long lookback on a long history) submitted last, as a sweep over a grid naturally submits its largest parameters. The work
adds up to 3\,649 core-hours on 384 cores, a lower bound of 9.50~hours.

\begin{omfigure}
\begin{tikzpicture}
\begin{axis}[width=12cm,height=5.6cm,xlabel={hours after submission},ylabel={busy cores (log scale)},xmin=0,xmax=21,
  ymode=log,ymin=0.7,ymax=600,unbounded coords=jump,
  grid=major,grid style={gray!20},tick label style={font=\small},label style={font=\small},table/col sep=comma]
\addplot[color=omThm,very thick,const plot] table[x=t_h,y=busy]{figdata/platforms/13-distributed-compute-and-schedulers/profile_fifo.csv};
\addplot[color=omDef,thick,densely dashed,const plot] table[x=t_h,y=busy]{figdata/platforms/13-distributed-compute-and-schedulers/profile_lpt.csv};
\node[omThm,font=\scriptsize,anchor=south] at (axis cs:14,45) {first in, first out};
\node[omDef,font=\scriptsize,anchor=north west,align=left] at (axis cs:10.6,330) {longest first, with backups:\\ends at 10.0};
\draw[black,dotted] (axis cs:9.50,0.7) -- (axis cs:9.50,600) node[pos=0.3,right,font=\scriptsize]{lower bound};
\end{axis}
\end{tikzpicture}

\omcaption{Busy cores over time for the same sweep under two schedules, sampled every ten minutes, on a log scale. First in, first out keeps
the cluster full for about nine hours, reaches the forty long tasks only after twelve, and runs them on an almost idle cluster until
the last ends at 19.8~hours; longest first with speculative backups ends at 10.0~hours. Data: \texttt{fig\_jobgraph.py}.}\label{fig:pl:distributed-compute-and-schedulers:profile}
\end{omfigure}

\Cref{tab:pl:distributed-compute-and-schedulers:policies} is the sweep under each policy. First in, first out reaches the last of the forty long
tasks only after twelve hours, and the sweep ends at 19.8~hours --- 475.0 node-hours of a cluster held for a job that needed 3\,649
core-hours. Longest first starts the long tasks at once and ends at 11.8~hours.
\omterm{def:pl:distributed-compute-and-schedulers:stealing}{Work stealing}, which balances load but does not reorder it, ends at 16.3.

\begin{table}[ht]
\centering\small
\begin{tabular}{lrrrrr}
\toprule
policy & \multicolumn{2}{c}{\omterm{def:pl:distributed-compute-and-schedulers:makespan}{makespan} (hours)} & \multicolumn{2}{c}{node-hours} & backups \\
 & plain & with backups & plain & with backups & \\
\midrule
first in, first out & 19.79 & 20.32 & 475.0 & 487.6 & 152 \\
longest first & 11.80 & 10.02 & 283.3 & 240.4 & 105 \\
\omterm{def:pl:distributed-compute-and-schedulers:stealing}{work stealing} & 16.33 & 16.40 & 391.9 & 393.5 & 145 \\
longest first, no slow node & 10.32 & --- & 247.6 & --- & --- \\
\midrule
lower bound & \multicolumn{2}{c}{9.50} & & & \\
\bottomrule
\end{tabular}
\caption{The 10\,000-task sweep on 384 cores under each policy, without and with speculative backups (a backup after 2.5 times the estimate),
and longest first on a cluster without the slow node. Node-hours are the cluster's 24 nodes held for the \omterm{def:pl:distributed-compute-and-schedulers:makespan}{makespan}.}
\label{tab:pl:distributed-compute-and-schedulers:policies}
\end{table}

\begin{omfigure}
\begin{tikzpicture}
\begin{axis}[width=12cm,height=5.4cm,ybar,bar width=12pt,ylabel={makespan (hours)},ymin=0,ymax=26,
  symbolic x coords={FIFO,LPT,work stealing},xtick=data,enlarge x limits=0.25,
  tick label style={font=\small},label style={font=\small},grid=major,grid style={gray!20},table/col sep=comma,
  legend style={font=\footnotesize,at={(0.5,1.03)},anchor=south,legend columns=2},nodes near coords,
  nodes near coords style={font=\scriptsize,rotate=90,anchor=west,/pgf/number format/fixed,/pgf/number format/fixed zerofill,/pgf/number format/precision=2},point meta=y,ybar=4pt]
\addplot[fill=omThm!45,draw=omThm] table[x=policy,y=plain]{figdata/platforms/13-distributed-compute-and-schedulers/makespans.csv};
\addlegendentry{plain}
\addplot[fill=omDef!45,draw=omDef] table[x=policy,y=backups]{figdata/platforms/13-distributed-compute-and-schedulers/makespans.csv};
\addlegendentry{with backups}
\draw[black,dashed] ({rel axis cs:0,0}|-{axis cs:FIFO,9.50}) -- ({rel axis cs:1,0}|-{axis cs:FIFO,9.50})
  node[pos=0.335,below,font=\tiny]{lower bound 9.50};
\end{axis}
\end{tikzpicture}

\omcaption{\omterm{def:pl:distributed-compute-and-schedulers:makespan}{Makespan} of the sweep by policy, without and with speculative backups (after 2.5 times the estimate), against the lower bound.
First in, first out (FIFO) and \omterm{def:pl:distributed-compute-and-schedulers:stealing}{work stealing} run the long tasks last; longest first (LPT) runs them first. Data:
\texttt{fig\_jobgraph.py}.}\label{fig:pl:distributed-compute-and-schedulers:makespans}
\end{omfigure}

\begin{definition}[Speculative execution]\label{def:pl:distributed-compute-and-schedulers:spec}
\emph{Speculative execution}\index{speculative execution} launches a second copy of a task that is running much longer than expected, on
another machine; the first copy to finish is kept and the other is killed.
\end{definition}

Google's MapReduce paper named the problem and the cure: a \omterm{def:pl:distributed-compute-and-schedulers:makespan}{straggler} can be a machine with a bad disk whose read rate falls from 30 to
1~MB/s, and when a job is close to completion the master schedules backup executions of the remaining in-progress tasks; their sort took
44\% longer with backups disabled. The threshold matters. The chapter's backups start when a task has run 2.5 times its estimate: with
estimates right to within $e^{0.3z}$, an honest task passes that point with probability 0.11\% ($z > 3.05$), so backups go mostly to tasks on
the slow node. Longest first gains 1.8~hours from them; first in, first out loses 0.5, because its long tasks are late to start, not
slow to run, and the backups only take cores; the wasted work --- killed copies, failed and retried attempts --- is 91 to 130
core-hours, 2.5 to 3.6\% of the sweep.

\omcode{code/firm/jobgraph/firm_jobgraph.py}{206}{227}{Speculative execution in the simulator: a late task's backup goes to a free core on
another node before the queue is served.}\label{lst:pl:distributed-compute-and-schedulers:spec}

\begin{example}[The sweep that finished on Saturday]\label{ex:pl:distributed-compute-and-schedulers:sweep}
Longest first with backups ends at 10.02~hours, 5.5\% above the lower bound of 9.50 hours. Without the slow node, but also without
backups, it ends at 10.32: the backups recover more than the slow node costs, because they also rescue the tasks whose estimates were
wrong and those that failed. First in, first out on the same cluster ends at 19.79~hours, twice as long, and holds 475 node-hours against
240. The cluster had been almost idle for three and a half hours: from 16.3~hours after submission, fewer than a tenth of its cores were
busy.
\end{example}

\section{Fair share between teams}

A second team submits 500 backtests an hour after the first team's sweep. Under first in, first out they wait behind all 10\,000: half of them
are done 8.28~hours after submission and the last 11.4~hours after. Under fair share with equal targets, every core that frees goes to the
second team until it holds half the cluster: half its tasks are done in 0.59~hours and all in 2.64, while the first team's sweep ends at
16.68~hours instead of 19.79 --- earlier, not later, because interleaving the second team's tasks changes where and when the first
team's long tasks run, and that, not the extra 5\% of work, sets its end.

\section{Caching, failure and cost}

When every backtest recomputes its features (twenty minutes each), the sweep needs 6\,922 core-hours and 19.3~hours; written as a \omterm{def:pl:distributed-compute-and-schedulers:graph}{task graph}
with a hundred feature stages it needs 3\,514 core-hours and 11.4~hours; the next sweep over the same data finds the hundred stages in the
cache and needs 3\,475 core-hours and 10.3~hours. Failures cost what was lost: a failed attempt is retried from its start, so a failure of a
seven-hour task costs up to seven hours, which is why long tasks write checkpoints (Book~12, chapter~23). The cost of a sweep is the capacity
it holds: on a cluster rented by the node-hour, the difference between first in, first out and longest first with backups is 235 node-hours for
the same answers.

\section{Tutorial: a sweep on a simulated cluster}\label{tut:pl:distributed-compute-and-schedulers:tut}

\begin{tutorial}
\textbf{Goal.} Calibrate a sweep on real backtests, schedule it under four policies on a simulated cluster, and measure what backups,
fair share and caching change. \textbf{End state:} \cref{tab:pl:distributed-compute-and-schedulers:policies} and
\cref{fig:pl:distributed-compute-and-schedulers:profile}.
\begin{tutsteps}
\item \textbf{Calibrate}: \texttt{bench\_jobgraph.py} runs 24 small backtests through \texttt{LocalExecutor} with one worker;
\texttt{measured\_sigma()}.
\item \textbf{Sweep}: \texttt{sweep(sigma)}, \texttt{cluster()}, \texttt{lower\_bound}.
\item \textbf{Policies}: \texttt{policies(sigma)}; \texttt{profile(sigma)} for the busy-core curves.
\item \textbf{Teams}: \texttt{fair\_share(sigma)}.
\item \textbf{Caching}: \texttt{caching(sigma)} on the \omterm{def:pl:distributed-compute-and-schedulers:graph}{task graph}.
\end{tutsteps}
\textbf{What to change next.} Give the scheduler perfect estimates and see how close longest first comes to the bound; give it no estimates
(every task the median) and see what remains of its advantage.
\end{tutorial}

\section{Build: the job graph}\label{bld:pl:distributed-compute-and-schedulers:jobgraph}

\begin{build}
\bfield{Purpose} A task-graph and scheduling library the platform uses to plan and run sweeps, and a simulator to choose policies before
spending node-hours.
\bfield{Interface} \texttt{Task}, \texttt{Cluster}, \texttt{simulate(tasks, cluster, policy, spec, cache) -> Schedule};
\texttt{FIFO}, \texttt{LPT}, \texttt{FairShare}, \texttt{WorkStealing}; \texttt{lower\_bound}; \texttt{utilisation};
\texttt{LocalExecutor}.
\bfield{Rules} The scheduler sees estimates, never durations; a failed or preempted attempt is retried from its start; a backup runs on
another node and the loser is killed and counted as waste; a cache hit takes no time; everything is seeded.
\bfield{Acceptance tests} \texttt{code/firm/jobgraph/tests/}: Graham's example ($3,3,2,2,2$ on two cores: 7 against an optimum of 6); first
in, first out against longest first; dependencies and the chain bound; fair share splitting cores; \omterm{def:pl:distributed-compute-and-schedulers:stealing}{work stealing} on an unbalanced deal;
backups rescuing a slow node; retries; cache hits; the local executor's order.
\bfield{Stretch} Memory classes and multi-core tasks with backfilling; preemptible (spot) capacity at a lower price; usage decaying over time in
fair share.
\end{build}

\begin{omsources}
\item R.~L. Graham, ``Bounds on multiprocessing timing anomalies'', \emph{SIAM Journal on Applied Mathematics} 17(2), 1969.
\item J.~Dean and S.~Ghemawat, ``MapReduce: simplified data processing on large clusters'', OSDI 2004.
\item J.~Dean and L.~A. Barroso, ``The tail at scale'', \emph{Communications of the ACM} 56(2), 2013.
\item R.~D. Blumofe and C.~E. Leiserson, ``Scheduling multithreaded computations by work stealing'', \emph{Journal of the ACM} 46(5), 1999.
\item Slurm documentation, ``Fair Tree Fairshare Algorithm''.
\end{omsources}

\section{Exercises}

\begin{exercise}[$\star$]\label{exo:pl:distributed-compute-and-schedulers:1}
Five tasks of durations 3, 3, 2, 2, 2 run on two cores. What is the \omterm{def:pl:distributed-compute-and-schedulers:makespan}{makespan} under longest first, and what is the best possible?
\end{exercise}

\begin{exercise}[$\star$]\label{exo:pl:distributed-compute-and-schedulers:2}
The sweep's median task takes fifteen minutes. How long do the forty long tasks take on a normal node, and on the slow node?
\end{exercise}

\begin{exercise}[$\star$]\label{exo:pl:distributed-compute-and-schedulers:3}
Why is the lower bound of the chapter's sweep the total work divided by the cores, and not the longest task?
\end{exercise}

\begin{exercise}[$\star\star$]\label{exo:pl:distributed-compute-and-schedulers:4}
Why does the backup threshold of 2.5 times the estimate rarely fire for a healthy task, and what would a threshold of 1.5 do?
\end{exercise}

\begin{exercise}[$\star\star$]\label{exo:pl:distributed-compute-and-schedulers:5}
Why does \omterm{def:pl:distributed-compute-and-schedulers:stealing}{work stealing} finish later than longest first, although no core is idle while work remains queued?
\end{exercise}

\begin{exercise}[$\star\star$]\label{exo:pl:distributed-compute-and-schedulers:6}
Fair share lets the second team finish sooner. What does it do to the first team's end, and why?
\end{exercise}

\begin{exercise}[$\star\star\star$]\label{exo:pl:distributed-compute-and-schedulers:7}
\emph{Coding.} Run \texttt{policies} with the long tasks submitted first instead of last. What happens to first in, first out, and to
longest first?
\end{exercise}

\begin{exercise}[$\star\star\star$]\label{exo:pl:distributed-compute-and-schedulers:8}
\emph{Find the flaw.} ``The sweep ran all night on a cluster of 384 cores and used only 3\,649 core-hours, so the cluster is too big: we
will halve it.''
\end{exercise}

\section{Problem: The Sweep That Finished on Saturday}

\begin{problem}[{Weekend problem --- ten thousand backtests and a slow node}]\label{pb:pl:distributed-compute-and-schedulers:1}
The chapter's sweep, cluster and policies.

\medskip\noindent\textbf{Part I --- The workload.}
\begin{enumerate}
\item Why is a \omterm{def:pl:distributed-compute-and-schedulers:sweep}{parameter sweep} \omterm{def:pl:distributed-compute-and-schedulers:sweep}{embarrassingly parallel}?
\item How was the spread of task durations obtained?
\item What is the total work, and the lower bound on 384 cores?
\item Where are the forty long tasks in the submission order, and why?
\item What makes a task a \omterm{def:pl:distributed-compute-and-schedulers:makespan}{straggler} here, in three ways?
\end{enumerate}

\medskip\noindent\textbf{Part II --- Policies.}
\begin{enumerate}[resume]
\item What does Graham's analysis guarantee for list scheduling and for longest first?
\item What does longest first need that first in, first out does not?
\item Why does \omterm{def:pl:distributed-compute-and-schedulers:stealing}{work stealing} not fix the long tasks?
\item What does fair share change between two teams?
\item When does \omterm{def:pl:distributed-compute-and-schedulers:spec}{speculative execution} help, and when does it waste?
\end{enumerate}

\medskip\noindent\textbf{Part III --- The numbers.}
\begin{enumerate}[resume]
\item Give the \omterm{def:pl:distributed-compute-and-schedulers:makespan}{makespan} of each policy with and without backups.
\item How long was the cluster almost idle under first in, first out?
\item How far from the lower bound is the best schedule, and why?
\item What does fair share do for the second team's first half of tasks?
\item What do the \omterm{def:pl:distributed-compute-and-schedulers:graph}{task graph} and the cache save?
\end{enumerate}

\medskip\noindent\textbf{Part IV --- The verdict.}
\begin{enumerate}[resume]
\item State the \emph{named result}: the \omterm{def:pl:distributed-compute-and-schedulers:makespan}{makespan} and node-hours of the sweep under each policy, the gain from \omterm{def:pl:distributed-compute-and-schedulers:spec}{speculative execution}, and
the distance of the best schedule from the lower bound.
\item What would you change in the platform's default scheduling of sweeps?
\item What would you ask users to provide with each sweep?
\item How would you decide the size of the cluster?
\item In one sentence: what decides when a sweep ends?
\end{enumerate}
\end{problem}

\section{Interview questions}

\begin{interviewq}[$\star$ \iqroles{developer}]\label{iq:pl:distributed-compute-and-schedulers:1}
Your sweep's last 1\% of tasks takes as long as the first 99\%. What do you do?
\end{interviewq}

\begin{interviewq}[$\star\star$ \iqroles{developer}]\label{iq:pl:distributed-compute-and-schedulers:2}
What is \omterm{def:pl:distributed-compute-and-schedulers:spec}{speculative execution}, and when is it a bad idea?
\end{interviewq}

\begin{interviewq}[$\star\star$ \iqroles{developer, researcher}]\label{iq:pl:distributed-compute-and-schedulers:3}
Two teams share a cluster; one submits huge sweeps. How do you keep the other productive?
\end{interviewq}

\begin{interviewq}[$\star\star$ \iqroles{researcher}]\label{iq:pl:distributed-compute-and-schedulers:4}
How would you structure ten thousand backtests that share expensive feature computations?
\end{interviewq}

\begin{interviewq}[$\star\star\star$ \iqroles{developer}]\label{iq:pl:distributed-compute-and-schedulers:5}
Design the scheduling layer of a research platform's \omterm{def:pl:distributed-compute-and-schedulers:cluster}{compute cluster}.
\end{interviewq}
